% Rapport technique AgriVoix
% Compilation : pdfLaTeX (recommandé, Overleaf par défaut) ou XeLaTeX.
% Mise en forme : Times New Roman 12 pt, interligne 1,5, sortie PDF 1.5.
\documentclass[12pt,a4paper]{article}

\usepackage{iftex}
\ifPDFTeX
  \pdfminorversion=5
  \usepackage[T1]{fontenc}
  \usepackage[utf8]{inputenc}
  \usepackage{newtxtext,newtxmath} % police Times
  \usepackage[safe]{tipa}          % caractères dioula : ɛ, ɔ, ɲ
  \DeclareUnicodeCharacter{025B}{\textepsilon}
  \DeclareUnicodeCharacter{0254}{\textopeno}
  \DeclareUnicodeCharacter{0272}{\textltailn}
\else
  \usepackage{fontspec}
  \IfFontExistsTF{Times New Roman}{\setmainfont{Times New Roman}}{\setmainfont{TeX Gyre Termes}}
\fi

\usepackage[french]{babel}
% babel-french met un tiret long devant chaque élément de liste : on le désactive.
\frenchsetup{StandardLists=true}
\usepackage[a4paper,margin=2.5cm]{geometry}
\usepackage{setspace}
\onehalfspacing
\usepackage{booktabs}
\usepackage{tabularx}
\usepackage{array}
\usepackage{xcolor}
\usepackage{tikz}
\usetikzlibrary{positioning,arrows.meta,babel}
\usepackage{enumitem}
\setlist{itemsep=2pt,topsep=4pt}
\setlist[itemize]{label=-} % puces : tiret simple (touche 6)
\usepackage{fancyvrb}
\usepackage[hidelinks]{hyperref}

\definecolor{agri}{HTML}{1F8A4C}
\newcolumntype{L}{>{\raggedright\arraybackslash}X}

\title{\textbf{AgriVoix}\\[4pt]\large Rapport technique\\[2pt]
\normalsize Plateforme vocale de mise en relation agricole en dioula, français et anglais}
\author{Équipe AgriVoix}
\date{Novembre 2026}

\begin{document}

\maketitle
\thispagestyle{empty}

\begin{abstract}
AgriVoix permet à des producteurs et acheteurs agricoles de Côte d'Ivoire de publier une annonce
de vente ou d'achat en parlant, en dioula, en français ou en anglais. L'application transcrit le
vocal, en comprend l'intention (vendre ou acheter, produit, quantité, lieu), met automatiquement en
relation l'offre et la demande, puis ouvre une conversation entre les deux personnes. Un assistant
agricole répond aussi aux questions posées à la voix ou par photo, dans la langue de l'utilisateur.
Ce rapport décrit l'architecture, la chaîne de traitement de la parole et, en particulier, la
spécialisation agricole des modèles Djelia à l'aide d'un dictionnaire agricole français-dioula
construit pour le projet, qui évite que le système se comporte comme un traducteur généraliste.
Il présente enfin la feuille de route : intégrer progressivement les langues locales africaines,
à commencer par le wolof dès décembre 2026, pour faire d'AgriVoix l'application de référence de la
vente agricole en Afrique.
\end{abstract}

\tableofcontents
\newpage

\section{Contexte et objectifs}

Une grande partie des petits producteurs vivriers d'Afrique de l'Ouest s'exprime d'abord dans une
langue nationale comme le dioula, et lit ou écrit peu le français. Les plateformes de vente en ligne
reposent pourtant sur des formulaires écrits. AgriVoix part du principe inverse : la voix est
l'interface principale.

Les objectifs fixés sont les suivants :
\begin{itemize}
  \item publier une annonce en maintenant un seul bouton, sans rien saisir ;
  \item comprendre le dioula, le français et l'anglais, sans que l'utilisateur ait à choisir sa langue ;
  \item reconnaître correctement le vocabulaire agricole local (noms de cultures, de produits) ;
  \item mettre en relation automatiquement un vendeur et un acheteur du même produit ;
  \item permettre aux deux parties d'échanger (appel, WhatsApp, messages écrits et vocaux) ;
  \item offrir un conseiller agricole accessible à l'oral et par photo ;
  \item ne rien simuler : comptes, audios, annonces et messages sont réellement stockés.
\end{itemize}

\section{Architecture générale}

Le système est découpé en trois parties indépendantes (figure~\ref{fig:archi}).

\begin{figure}[h]
\centering
\begin{tikzpicture}[
  node distance=1.1cm and 1.4cm,
  box/.style={draw=agri, rounded corners, thick, align=center, minimum width=3.6cm, minimum height=1.1cm, font=\small},
  ext/.style={draw=gray, rounded corners, align=center, minimum width=3.2cm, minimum height=0.9cm, font=\small},
  arr/.style={-{Stealth}, thick}
]
\node[box] (front) {Application web\\Next.js (Vercel)};
\node[box, right=of front] (back) {API Python\\FastAPI (Railway)};
\node[box, right=of back] (db) {Supabase\\PostgreSQL + Stockage};
\node[ext, below=of front] (whisper) {Whisper (local)\\fr / en};
\node[ext, below=of back] (djelia) {Djelia\\ASR, traduction, voix};
\node[ext, below=of db] (claude) {Claude\\analyse, conseils};
\draw[arr] (front) -- node[above, font=\scriptsize]{HTTPS} (back);
\draw[arr] (back) -- node[above, font=\scriptsize]{RPC} (db);
\draw[arr] (back) -- (whisper);
\draw[arr] (back) -- (djelia);
\draw[arr] (back) -- (claude);
\end{tikzpicture}
\caption{Architecture d'AgriVoix}
\label{fig:archi}
\end{figure}

\begin{description}
  \item[Application web.] Next.js 16 et React, pensée pour le téléphone, déployée sur Vercel.
  Le navigateur n'appelle qu'une route relais \texttt{/api/*} du même domaine, qui transmet les
  requêtes au serveur Python. Aucune clé d'API n'est exposée côté client.
  \item[Serveur Python.] FastAPI et uvicorn, déployé sur Railway avec ffmpeg. Il orchestre la chaîne
  vocale, les appels aux modèles et l'accès à la base de données.
  \item[Base de données.] Supabase (PostgreSQL et stockage de fichiers) pour les profils, sessions,
  annonces, messages et fichiers audio.
\end{description}

Les services externes sont Djelia (reconnaissance vocale, traduction et synthèse vocale en dioula),
l'API Claude d'Anthropic (modèle \texttt{claude-haiku-5-5}, analyse des annonces et conseils) et
OpenStreetMap Nominatim (nom de la localité à partir des coordonnées GPS). La transcription du
français et de l'anglais est faite sur le serveur lui-même par Whisper, sans service tiers.

\section{Chaîne de traitement de la parole}

\subsection{Capture et préparation de l'audio}

L'enregistrement suit le principe du talkie-walkie : l'utilisateur maintient le bouton micro, parle,
puis relâche, et le vocal part automatiquement. Les enregistrements de moins de 0,7 seconde sont
ignorés (appui accidentel). Le navigateur produit selon les appareils du WebM, de l'OGG ou du MP4.

À la réception, un seul appel à ffmpeg produit deux fichiers :
\begin{itemize}
  \item un WAV mono à 16 kHz, format attendu par les modèles de reconnaissance vocale ;
  \item un MP3 mono à 64 kbit/s, conservé dans le stockage Supabase pour la réécoute.
\end{itemize}

\subsection{Détection automatique de la langue}

L'utilisateur ne choisit pas sa langue : elle est détectée sur l'audio. Le serveur utilise
faster-whisper (modèle \texttt{small}, exécuté en local sur le processeur, quantifié en int8).
Whisper reconnaît le français et l'anglais mais ne connaît pas le dioula. La règle de décision
exploite cette propriété :
\begin{itemize}
  \item si la probabilité du français ou de l'anglais dépasse 0,6, le vocal est traité dans cette langue ;
  \item sinon, il est considéré comme du dioula et confié à Djelia.
\end{itemize}
Le seuil est réglable par variable d'environnement. Pour les questions écrites, une détection
lexicale simple distingue le français de l'anglais.

\subsection{Parcours dioula}

\begin{enumerate}
  \item Transcription du vocal par le modèle de reconnaissance vocale Djelia \texttt{sunjata-1}.
  \item Normalisation du texte dioula à l'aide du dictionnaire agricole (section~\ref{sec:dico}).
  \item Traduction dioula vers français par le modèle Djelia \texttt{banjugu-1}.
  \item Analyse ou réponse par Claude, avec le glossaire agricole dans ses instructions.
  \item Pour l'assistant : traduction de la réponse vers le dioula, puis synthèse vocale par le
  modèle Djelia \texttt{jifili-1} (voix \og moussa \fg{}).
\end{enumerate}

\subsection{Parcours français et anglais}

Le vocal est transcrit par Whisper. Pour fiabiliser les mots courts, le modèle reçoit un contexte
agricole (liste de produits, unités) : sans lui, \og du riz \fg{} était transcrit \og durée \fg{}.
Claude répond dans la langue de l'utilisateur, et la réponse est lue par la synthèse vocale native
du navigateur (Web Speech API), sans coût serveur.

\section{Spécialisation agricole de Djelia}
\label{sec:dico}

\subsection{Problème : un modèle généraliste}

Les modèles Djelia sont entraînés pour le dioula et le bambara en général. Utilisés tels quels sur
des annonces agricoles, ils commettent des erreurs qui rendent l'annonce inutilisable. Le cas le plus
fréquent observé pendant les tests est le mot \textit{kaba} (maïs), traduit selon les phrases par
\og tabac \fg{}, \og cacao \fg{}, \og chou \fg{} ou \og canne à sucre \fg{}. Une erreur de ce type
fausse directement la mise en relation : un vendeur de maïs serait présenté à un acheteur de cacao.

\subsection{Solution : un dictionnaire agricole personnalisé}

Nous avons construit un dictionnaire agricole français-dioula propre au projet, ciblé sur la
Côte d'Ivoire et établi à partir d'un lexique dioula-français. Il recense 82 termes, dont les
principales cultures, et 39 variantes orthographiques ou régionales rattachées à leur forme de
référence (tableau~\ref{tab:dico}).

\begin{table}[h]
\centering
\begin{tabular}{lc}
\toprule
Catégorie & Nombre d'entrées \\
\midrule
Céréales & 8 \\
Tubercules & 6 \\
Légumineuses & 6 \\
Légumes & 18 \\
Fruits & 17 \\
Cultures industrielles (cacao, café, anacarde\dots) & 9 \\
Plantes aromatiques & 4 \\
Vocabulaire agricole (champ, récolte, feuille\dots) & 6 \\
Couleurs (description des maladies, de la maturité) & 8 \\
\midrule
\textbf{Total} & \textbf{82} \\
\bottomrule
\end{tabular}
\caption{Composition du dictionnaire agricole}
\label{tab:dico}
\end{table}

Chaque entrée associe le terme français, sa forme dioula de référence et ses variantes, par exemple :
\begin{itemize}
  \item Maïs : \textit{kaba}, variantes \textit{mɔnsɔ}, \textit{mɔnɔnɔ}, \textit{maanɔ} ;
  \item Riz : \textit{malo}, variante \textit{kini} ;
  \item Sorgho : \textit{binbiri}, variante \textit{kɛnde}.
\end{itemize}

\subsection{Intégration à trois niveaux de la chaîne}

Le dictionnaire spécialise la chaîne Djelia sur le domaine agricole à trois endroits.

\begin{enumerate}
  \item \textbf{Avant la traduction : normalisation.} Chaque variante reconnue dans la transcription
  Djelia est remplacée par la forme de référence, avec une recherche par mots entiers, insensible à
  la casse, qui traite d'abord les expressions les plus longues. Le modèle de traduction reçoit ainsi
  toujours le terme qu'il connaît le mieux.
  \item \textbf{Pendant la génération : glossaire injecté.} La liste complète des correspondances
  français-dioula est ajoutée aux instructions de Claude, avec la consigne de ne jamais confondre
  une plante avec un plat cuisiné. Les réponses de l'assistant, puis leur traduction en dioula,
  emploient donc les bons noms de cultures.
  \item \textbf{Après la traduction : correction de l'annonce.} Pour une annonce vocale, Claude reçoit
  à la fois la transcription dioula et la traduction automatique, avec la consigne que la
  transcription dioula et le glossaire font foi en cas de désaccord. Il produit une version française
  corrigée et les champs structurés de l'annonce.
\end{enumerate}

\subsection{Résultat}

Sur la phrase de test \og N bɛ kaba kilo kɛmɛ duuru feere \fg{} (je vends cinq cents kilos de maïs),
la traduction automatique seule proposait un autre produit. Avec la chaîne spécialisée, l'annonce
est correctement classée en vente de maïs. Le dictionnaire transforme ainsi un traducteur généraliste
en outil adapté au vocabulaire des marchés agricoles, sans dépendre de la qualité de la traduction
mot à mot.

Ce dispositif est également facile à enrichir : ajouter une culture ou une variante régionale
consiste à modifier un fichier JSON, sans réentraîner aucun modèle.

\section{Analyse des annonces et mise en relation}

\subsection{Extraction structurée}

Chaque annonce, vocale ou écrite, est analysée par Claude avec une sortie contrainte par un schéma
JSON, ce qui garantit un résultat exploitable sans traitement de texte fragile :

\begin{Verbatim}[fontsize=\small]
{ "kind":     "offer" | "need",      vente ou recherche
  "product":  Maïs, Manioc, Riz, Tomate, Igname, Banane,
              Cacao, Café ou Autre,
  "quantity": "500 kg",
  "place":    "Korhogo",
  "french":   version française corrigée de l'annonce }
\end{Verbatim}

Une annonce en anglais est traduite en français dans le champ \texttt{french}, afin que tous les
utilisateurs puissent la lire et qu'elle participe à la même mise en relation.

\subsection{Algorithme de mise en relation}

La mise en relation est exécutée dans la base de données, dans la même transaction que
l'enregistrement de l'annonce. Une annonce en attente est associée à la plus ancienne annonce en
attente qui porte sur le même produit, avec l'intention opposée (vente face à achat) et publiée par
un autre utilisateur. Les annonces de catégorie \og Autre \fg{} ne sont pas associées
automatiquement.

Les lignes concernées sont verrouillées (\texttt{FOR UPDATE SKIP LOCKED}) : si deux annonces arrivent
en même temps, une même offre ne peut pas être attribuée à deux acheteurs. Lorsqu'une annonce est
annulée, l'annonce associée repasse en attente et peut trouver un autre partenaire.

\subsection{Échanges entre les deux parties}

Une fois l'association faite, chacun voit la fiche de l'autre (nom, téléphone, localité, distance)
avec trois actions : appeler, écrire sur WhatsApp ou ouvrir la conversation intégrée. Celle-ci
accepte des messages écrits et vocaux ; les vocaux sont stockés en MP3 dans Supabase.

\section{Données et sécurité}

\subsection{Modèle de données}

\begin{table}[h]
\centering
\begin{tabularx}{\textwidth}{lL}
\toprule
Table & Contenu \\
\midrule
\texttt{profiles} & nom, téléphone (format international, unique), localité, coordonnées GPS, code de connexion haché \\
\texttt{sessions} & une session par appareil connecté (empreinte du jeton) \\
\texttt{annonces} & intention, produit, quantité, lieu, transcriptions, chemin de l'audio, statut, annonce associée \\
\texttt{messages} & messages écrits et vocaux entre deux annonces associées \\
\texttt{assistant\_vocaux} & questions vocales posées à l'assistant, transcriptions et réponses \\
\bottomrule
\end{tabularx}
\caption{Tables principales}
\end{table}

\subsection{Mesures de sécurité}

\begin{itemize}
  \item \textbf{Aucun accès direct à la base depuis le navigateur.} La sécurité au niveau des lignes
  (RLS) est activée sans aucune règle d'accès public. Toutes les opérations passent par des fonctions
  SQL exécutables uniquement par le serveur.
  \item \textbf{Sessions.} À l'inscription ou à la connexion, le serveur génère un jeton aléatoire de
  256 bits. Seule son empreinte SHA-256 est stockée. Chaque appareil a sa propre session, révoquée
  à la déconnexion.
  \item \textbf{Connexion par code.} En l'absence de vérification par SMS, un code à 4 chiffres est
  demandé à l'inscription. Il est haché avec bcrypt dans la base. Après 5 erreurs, le compte est
  bloqué 15 minutes, et le message d'erreur est le même que le numéro existe ou non.
  \item \textbf{Audios privés.} Le stockage des fichiers est privé. La lecture se fait par des liens
  signés valables 6 heures, réservés aux personnes concernées.
  \item \textbf{Clés d'API.} Les clés Djelia, Anthropic et Supabase restent sur le serveur Python.
\end{itemize}

\section{Interface utilisateur}

L'application compte trois onglets, dans un style volontairement sobre inspiré d'iOS :
\begin{description}
  \item[Vocal.] Un grand bouton micro au centre, à maintenir pour parler, et une saisie écrite
  facultative.
  \item[Mes vocaux.] Les annonces publiées, leur état (en attente ou associée) et le contact du
  partenaire, avec une pastille pour les nouveautés.
  \item[Assistant.] Une conversation avec le conseiller agricole, par texte, photo ou voix, avec
  lecture des réponses à voix haute.
\end{description}

L'inscription demande le nom, le téléphone et le code ; la position est détectée automatiquement
par le GPS du téléphone, avec une saisie manuelle de la ville en cas de refus. Un bouton en haut à
droite donne accès au compte et à la déconnexion. Les nouveautés sont récupérées par interrogation
régulière du serveur (toutes les 10 secondes pour la liste, toutes les 4 secondes dans une
conversation).

\section{Déploiement}

\begin{itemize}
  \item \textbf{Application web} : Vercel, déploiement automatique à chaque mise à jour du dépôt Git.
  \item \textbf{Serveur Python} : Railway, construit avec Railpack (Python 3.12, ffmpeg installé). Le
  modèle Whisper est téléchargé au premier démarrage puis chargé en mémoire.
  \item \textbf{Base de données} : Supabase, schéma versionné sous forme de migrations SQL.
\end{itemize}

\section{Tests et validation}

\subsection{Parcours complet}

Un test automatisé de bout en bout exerce le système réel (serveur, modèles, base de données) :
inscription avec géolocalisation, refus d'un numéro déjà inscrit et d'un jeton invalide, annonce
vocale en dioula avec stockage et relecture de l'audio, mise en relation avec une annonce écrite,
messages écrits et vocaux, contrôle d'accès à la conversation, archivage d'une question vocale et
annulation. Les 16 vérifications sont réussies. Un second test couvre l'authentification
(inscription avec code, mauvais code, deux appareils, déconnexion, blocage) : 11 vérifications
réussies sur 11.

\subsection{Détection de la langue}

\begin{table}[h]
\centering
\begin{tabularx}{\textwidth}{lccL}
\toprule
Vocal & Langue détectée & Probabilité fr/en & Résultat \\
\midrule
Français & français & 0,99 & transcription correcte, réponse en français \\
Anglais & anglais & 1,00 & transcription correcte, réponse en anglais \\
Dioula & dioula & 0,32 & confié à Djelia, réponse vocale en dioula \\
\bottomrule
\end{tabularx}
\caption{Détection de la langue avec Whisper \texttt{small} (seuil 0,6)}
\end{table}

Une annonce écrite en anglais (\og I want to buy two tons of rice near Korhogo \fg{}) est correctement
analysée : achat, riz, 2 tonnes, Korhogo, avec la version française \og Je voudrais acheter deux
tonnes de riz près de Korhogo \fg{}.

Ces mesures ont été faites sur des voix de synthèse ; une campagne de tests avec de vrais locuteurs,
en particulier en dioula mêlé de mots français, reste à mener.

\section{Feuille de route : toutes les langues locales africaines}
\label{sec:roadmap}

\subsection{Ambition}

Le dioula n'est qu'un point de départ. L'objectif d'AgriVoix est de parler toutes les langues
dans lesquelles se négocient les récoltes en Afrique, et de devenir l'application numéro un de la
vente agricole sur le continent. Chaque nouvelle langue ouvre l'application à des producteurs qui
n'ont aujourd'hui accès à aucun outil numérique dans leur langue.

\subsection{Calendrier}

\begin{table}[h]
\centering
\begin{tabularx}{\textwidth}{lL}
\toprule
Échéance & Étape \\
\midrule
Novembre 2026 & Dioula, français et anglais en service (version actuelle) \\
Décembre 2026 & Intégration du wolof, lancement au Sénégal \\
2027 & Extension aux grandes langues agricoles d'Afrique de l'Ouest, puis du reste du continent, par ordre de nombre de producteurs \\
\bottomrule
\end{tabularx}
\caption{Calendrier d'intégration des langues}
\end{table}

\subsection{Une architecture prévue pour ajouter des langues}

La chaîne actuelle sépare déjà les rôles. Ajouter une langue revient à fournir quatre éléments,
sans modifier le reste de l'application (mise en relation, conversations, interface) :
\begin{enumerate}
  \item \textbf{un moteur de reconnaissance vocale} pour la langue ;
  \item \textbf{un moteur de traduction} vers le français, et retour ;
  \item \textbf{une voix de synthèse} pour lire les réponses de l'assistant ;
  \item \textbf{un dictionnaire agricole} dans le même format que le dictionnaire dioula, avec les
  noms locaux des cultures et leurs variantes.
\end{enumerate}
Le dictionnaire reste la pièce maîtresse : comme pour le dioula, c'est lui qui spécialise des
modèles généralistes sur le vocabulaire des marchés agricoles de chaque pays.

\subsection{Cas du wolof}

Pour le wolof, prévu en décembre 2026, les travaux sont les suivants :
\begin{itemize}
  \item constituer le dictionnaire agricole français-wolof (céréales, arachide, oignon, mil, niébé,
  maraîchage, unités de mesure locales) avec ses variantes ;
  \item retenir les moteurs de reconnaissance, de traduction et de synthèse vocale wolof. Les modèles
  multilingues ouverts couvrant le wolof, comme la famille Meta MMS, font partie des pistes étudiées ;
  \item adapter la liste des produits suivis aux filières sénégalaises.
\end{itemize}

\subsection{Faire évoluer la détection de la langue}

La règle actuelle (\og tout ce qui n'est ni du français ni de l'anglais est du dioula \fg{})
fonctionne avec une seule langue locale. Avec plusieurs langues locales, elle sera remplacée par un
modèle d'identification de la langue parlée couvrant les langues africaines, complété par la langue
préférée et le pays de l'utilisateur pour lever les ambiguïtés entre langues proches.

\section{Limites et perspectives}

\begin{itemize}
  \item \textbf{Reconnaissance vocale dioula.} Le modèle Djelia entend parfois \og ka ba \fg{} au lieu
  de \textit{kaba}. Le dictionnaire corrige les variantes connues, mais pas toutes les erreurs de
  segmentation.
  \item \textbf{Dioula mêlé de français.} Un vocal très mélangé peut être pris pour du français ; le
  seuil de détection devra être ajusté sur des enregistrements réels.
  \item \textbf{Vérification du numéro.} Il n'y a pas encore de vérification par SMS.
  \item \textbf{Temps réel.} Les nouveautés arrivent par interrogation régulière et non par
  notification instantanée.
  \item \textbf{Mise en relation.} L'association se fait par produit et par ordre d'arrivée ; la
  distance, la quantité et le prix pourraient être pris en compte.
  \item \textbf{Vers un entraînement dédié.} Les questions vocales archivées, avec leur transcription
  et leur traduction, constituent un corpus agricole réel. Une fois relu et validé, il pourra servir
  à entraîner spécifiquement les modèles de reconnaissance et de traduction sur le domaine agricole,
  en complément du dictionnaire.
\end{itemize}

\section{Conclusion}

AgriVoix montre qu'une plateforme agricole peut se passer de formulaires : un vocal suffit pour
publier une annonce, être mis en relation et obtenir un conseil, en dioula, en français ou en anglais.
Le point clé du système est la spécialisation agricole de la chaîne Djelia par un dictionnaire
personnalisé, appliqué avant la traduction, pendant la génération et lors de la correction des
annonces : elle rend la compréhension fiable sur le vocabulaire qui compte le plus pour la mise en
relation, les noms des produits.

Cette approche est reproductible : avec le wolof dès décembre 2026, puis les autres langues locales
africaines, AgriVoix vise à devenir l'application numéro un de la vente agricole en Afrique, celle
où chaque producteur peut vendre sa récolte en parlant sa propre langue.

\end{document}
